% Chapter 17 — Scope. Author: KLEPPMANN. Budget 3 pp (G4: +1 pp for the open-problems index).
% Discharges the Constitution v1.2 "Scope" section.
% Hosts the open-problems index and the constitutional-parking rule (G4).
% Target of Chapter~\ref{ch:requirements} V1 (model retention outside scope) via \label{ch:scope}.
% No thread episodes: a scope chapter cites; it does not re-run threads.
\section{Scope}\label{ch:scope}

This chapter discharges constitution clauses C-Scope.1--C-Scope.11 (the boundary of the specification; C-Scope.9 jointly with Chapter~\ref{ch:settlement}), the Authority clauses C-Auth.1--C-Auth.4 (C-Auth.1 and C-Auth.4 jointly with Chapter~\ref{ch:objective}), and C-6.6 (managed-account consequences, delegated by its own terms to the Managed-Account Companion, register line E75).

Scope fixes the boundary: what the specification covers, and what it reconciles against
at its perimeter but never performs. The constitution states the in-scope object in one
sentence --- the recording, valuation, and lifecycle of trading-book positions held at
fair value, and the interface that projects committed activity into settlement; everything
below is that sentence made exact, with the external authorities that lie outside it.

\subsection{In scope}

The specification covers the following components, each defined and proved in the chapter
named. Nothing in scope lives outside a chapter; nothing below is left to an implementation
to invent.

\begin{itemize}[leftmargin=*, itemsep=2.5pt]
\item \textbf{The map-then-fold architecture and the projection machinery} that derives
  every view from the log --- Chapter~\ref{ch:picture}.
\item \textbf{The immutable log and the objects} --- units, wallets, balances, moves,
  transactions, the coordinate vector --- Chapter~\ref{ch:objects}.
\item \textbf{The three machines} --- the Event Monitor, the Events Executor, and the
  Transaction Executor --- and the single-writer door --- Chapter~\ref{ch:machines}.
\item \textbf{The smart-contract framework}, its event-triggering discipline, and the
  purity and idempotence requirements --- Chapter~\ref{ch:contracts}.
\item \textbf{The three homes} --- ProductTerms, UnitStatus, PositionState --- and the
  conditions that render illegal states unrepresentable --- Chapter~\ref{ch:homes}.
\item \textbf{Valuation and profit and loss} --- NAV as a projection, path-independent PnL,
  and deposit neutrality --- Chapter~\ref{ch:valuation}.
\item \textbf{Ingestion, corporate actions, and the market data operator} --- the recorded
  observation, the registry, and the operator per (data-kind, event-kind) pair ---
  Chapter~\ref{ch:marketdata}.
\item \textbf{Collateral, margin, and lending} --- the regime key, the coverage invariant,
  and the determination/payment split --- Chapter~\ref{ch:collateral}.
\item \textbf{Virtual ledgers, strategies, and the total return swap} ---
  Chapter~\ref{ch:virtual}.
\item \textbf{The settlement interface} --- the projection of committed activity into
  settlement --- Chapter~\ref{ch:settlement}.
\item \textbf{Presentation and reporting projections} --- every reporting figure read from
  coordinates, units, and declared terms alone --- Chapter~\ref{ch:reporting}.
\item \textbf{Alignment with the Common Domain Model} --- shown, never adopted as authority
  --- Chapter~\ref{ch:cdm}.
\item \textbf{The invariant catalogue} and its executable properties ---
  Chapter~\ref{ch:invariants}.
\item \textbf{Testability and the executable-check regime} ---
  Chapter~\ref{ch:testability}.
\item \textbf{The minimum requirements} any conformant implementation must meet ---
  Chapter~\ref{ch:requirements}.
\end{itemize}

\subsection{Out of scope}

The following are external authorities. The ledger reconciles against each at its boundary;
it performs none of them. Each is a function of some other system, and the ledger's
closure holds precisely because it does not absorb them.

\begin{itemize}[leftmargin=*, itemsep=2.5pt]
\item \textbf{Price formation} --- the ledger records the outputs of pricing models and
  never runs one (Chapter~\ref{ch:valuation}; Chapter~\ref{ch:requirements}, V1); price
  discovery is external.
\item \textbf{Legal agreements} --- every right and obligation a contract creates is a
  unit (Chapter~\ref{ch:contracts}); the instrument's legal standing is decided outside
  the boundary.
\item \textbf{The settlement infrastructure itself} --- finality, payment-system access,
  custodian and CSD connectivity: the ledger says \emph{what} settles; the infrastructure
  decides how, and whether (Chapter~\ref{ch:settlement}).
\item \textbf{Regulatory submission} --- every report is a projection of the record
  (Chapter~\ref{ch:reporting}); its transmission and per-regime format are external.
\item \textbf{Reference-data authority} --- instrument, counterparty, and calendar
  mastering: reconciled against at the boundary, never performed.
\item \textbf{Model governance and retention} --- validation, approval, retention of
  models, and accounting policy. Reproducing a valuation from recorded inputs is
  unconditional; reproducing a model's number requires the model, whose retention is
  governance outside scope (Chapter~\ref{ch:requirements}, V1).
\end{itemize}

Five companion documents are named in the Exclusions Register as out-of-scope
deliverables --- expansions the specification points to but does not contain, and not
functions the ledger performs: the \emph{Worked-Examples Volume} (register lines E71 and
E72), holding the dividend-forecast, special-dividend, spin-off, and elective-merger
worked examples; the \emph{SBL Operations Companion} (register line E73), holding the full
securities-lending operational lifecycle --- locates, buy-ins, SFTR and CSDR discipline,
cascade recalls; the \emph{Track B graphs-interpreter companion} (register line E74),
holding the generic term-sheet interpreter, the CDM round-trip, and term-sheet rendering;
the \emph{Managed-Account Companion} (register line E75), holding managed-account fee
accrual and NAV attribution (the C-6.6 delegation); and the \emph{Simulation Companion}
(register line E76), holding the stochastic models, measures, calibration, and
variance-reduction techniques that drive simulated paths --- what is \emph{in} this
specification is the boundary that companion plugs into: branch construction, seed
discipline, result provenance (Chapter~\ref{ch:virtual}, Section~\ref{sec:simulated}).
Each is named so its absence is deliberate and auditable, never silent.

\subsection{The open-problems index and the constitutional-parking rule}

Two registers live here: open design questions, and parked constitutional conflicts.

\paragraph{The open-problems index.} The constitution requires that implementation
questions be named and left, not silently dropped; each is carried in full in the
\emph{Open-Problems and Escalation Register} companion. The top open risk is close-out and
netting algebra: per-netting-set claim units are first-class
(Chapter~\ref{ch:collateral}) precisely so close-out is representable, yet the algebra
that values each claim at close-out, nets per netting set, and decouples recovery from the
figure is not yet designed --- named here so it cannot be improvised. Named alongside it: the \textbf{right-of-use gate on collateral re-use} --- C-8.8 forbids
re-use of no-right-of-use collateral, yet the coverage net
(Chapter~\ref{ch:invariants}) bounds re-posting by possession alone, so a segregated
inflow is representable as deliverable; \emph{live and blocking} --- no implementation
enables re-use until the gate is built; and the \textbf{quantity-aware
settlement-obligation unit} --- a partial settlement leaves a settled quantity and a failed
residual the single lifecycle node cannot yet split into legs. Named further: the correction
algebra beyond the single compensating transaction; inter-ledger federation and
eliminations; \textbf{behavioural-event generation} --- a generator reaching the events
that are \emph{decisions} (default, elective exercise, merger election) must model the
deciding party, not just the market; \textbf{corporate-action generation} --- a simulated
corporate action is a generated \emph{terms version} crossing as a notice through the
term-adjustment operator, and its well-typed generator is open; \textbf{nested simulation}
--- a path that itself branches a simulated ledger (XVA off each exposure state), whose
seed discipline and provenance must compose across levels; and the remaining items of the
companion register. Naming
is the whole obligation; each is stated in the open at the chapter it touches, so none is a
gap hidden inside what the specification does claim.

\paragraph{The constitutional-parking rule.} The second register governs conflict with
Constitution v1.2 itself.

\begin{principle}[Constitutional parking]
A genuine conflict between this specification and the constitution is parked in this index
with the exact proposed amendment text --- never silently amended, never fudged. Amendment
is the owner's alone (the constitution's \emph{Authority and Amendment} section); parking
makes the conflict visible and hands the owner a precise, decidable choice.
\end{principle}

\emph{Current status: five closed, one ratified awaiting adoption.} Five conflicts parked
against v1.1 were resolved by the owner's v1.2; the Event Universe review parked one more,
ratified 2026-07-14. Each closure keeps its resolution --- an index that empties without a
trace would be indistinguishable from one never used:

\begin{itemize}
\item \textbf{PARK-1 (structural) --- CLOSED in v1.2.} Clause addressing: v1.2 embeds the
  margin identifiers, append-only; the separate rendering retired; every chapter opens with
  the clauses it discharges.
\item \textbf{PARK-2 (\S12) --- CLOSED in v1.2.} One timestamp where two axes were needed:
  C-12.1 restated --- both honest answers, indexing delegated here; the two-axis machinery
  retained (Theorem~\ref{thm:timetravel}).
\item \textbf{PARK-3 (\S4) --- CLOSED in v1.2.} Fixed schema against the netting collapse:
  C-4.10 restated --- ownership first, three admission tests, schema delegated here;
  Section~\ref{sec:coordinate-vector} demonstrates compliance
  (Theorem~\ref{thm:conservation}).
\item \textbf{PARK-4 (\S8) --- CLOSED in v1.2.} Unconditional deposit-neutrality against
  the day-one basis: C-8.7 adopted, conditional on fair value (Chapter~\ref{ch:valuation}).
\item \textbf{C-4.8 clarification --- CLOSED in v1.2.} The second door: every guaranteed
  observation a moveless transaction through the one Executor; the record is the log
  (Chapters~\ref{ch:objects}, \ref{ch:marketdata}, \ref{ch:invariants}).
\item \textbf{C-4.12 (the event universe) --- RATIFIED 2026-07-14, awaiting adoption.} No
  closure, no total routing, four loss paths on the captured side: event kinds now finite,
  closed, declared; routing total at registration; an undeclared event never processed and
  never lost --- quarantined through the one door. Discharged at
  Chapters~\ref{ch:marketdata}, \ref{ch:machines}, \ref{ch:testability},
  \ref{ch:requirements} (M8); adoption is the owner's sole act.
\end{itemize}

\noindent The index is live in its history: six conflicts parked with exact text, six
decided by the owner --- five adopted in v1.2, one ratified awaiting adoption --- the
resolutions above the evidence. A future conflict parks here the same way.

\subsection{Authority}

The direction of authority is one-way. This specification is rated against the
constitution, and against nothing else; the constitution is rated against nothing. Where a
future improvement requires departing from the constitution, the constitution is amended
first --- through the parking rule above and the owner's decision --- never the
specification bent in silence to fit it, and never the constitution bent to fit it.
